\documentclass[../main.tex]{subfiles}
\begin{document}

\section{Setup}
\label{sec:setup}

\subsection{Network}

We study a feedforward network with $L$ hidden layers whose preactivations are denoted by $\{ h^{(l)\mu}_{i} \}$, and whose outputs are denoted by $\{ f^{\mu}_\alpha \}$. To each neuron with index $i$ we associate a position in space given by $\mathbf{x}_i \in \mathbb{R}^d$. We will say more about the nature of this embedding later. Similarly, to each input index $a \in 1, \dots, m$ and output index $\alpha \in 1, \dots, n$, we associate a location in space, $\mathbf{x}_a$ and $\mathbf{x}_{\alpha}$, respectively.
\begin{align}
  h^{(1)\mu}_{i} =
    \frac{1}{\sqrt{m}} \sum_{\typo{a}=1}^{m} W^{(0)}_{i,a} x^{\mu}_a + \eta^{\mu}_i, \hfill
  &&
  h^{(l+1)\mu}_{i} =
    \frac{1}{\sqrt{N}} \sum_{\typo{j}=1}^N W^{(l)}_{i,j} \phi(h^{(l)\mu}_{j}), \hfill
  &&
  f^{\mu}_{\alpha} =
    \frac{1}{N} \sum_{\typo{j}=1}^{N} W^{(L)}_{\typo{\alpha},j} \phi(h^{(L)\mu}_{j})
  \label{eq:ff-network}
\end{align}
where $\eta$ is Gaussian noise, $\eta \sim \mathcal{N}(0, \sigma^2)$ (which we introduce to ensure positive-definiteness later on). We assume the network has trained to convergence under SGLD:
\begin{align*}
  \mathcal L[\Theta] &=
    \frac{N}{2 P n} 
    \sum_{\mu=1}^{P} \sum_{\alpha=1}^{n} 
      (f^{\mu}_{\alpha} - y^{\mu}_{\alpha})^2 
    + \frac{1}{\beta}  \mathcal{P}_{\text{weight}}(\Theta)
    + \frac{1}{\beta}  \mathcal{P}_{\text{rates}}(\Phi)
  \\
  \partial_s \Theta(s) &=
    -\nabla \mathcal{L}(\Theta(s))  + \sqrt{\frac{2}{\beta}} \eta(s),
  \quad \langle \eta_i(s) \eta_j(s') \rangle = \delta(s - s') \delta_{i,j}
\end{align*}
Such that we have the equilibrium distribution over parameters given by
\begin{align}
  p(\Theta) \xrightarrow{s \rightarrow \infty}  \frac{1}{\mathcal Z} e^{-\beta \mathcal{L}[\Theta]} \label{eq:theta-distr}
\end{align}
The $\mathcal{P}$ terms are regularisers which enact the spatial embedding:
\begin{align}
  \mathcal{P}_{\text{weight}} &=
    \frac12 \sum_{l=0}^{L} \sum_{i,j=1}^{\typo{\text{Post}(l)}} \sum_{a,b=1}^{\text{Pre}(l)} W^{(l)}_{i,a} M^{(l)}_{i,a,j,b} W^{(l)}_{j,b} \label{eq:prior-weights}
\end{align}
with 
\begin{align*}
  \text{Pre}(l) = \begin{cases}
    m & l = 0 \\
    N & \text{otherwise}
  \end{cases}
  &&
  \text{Post}(l) = \begin{cases}
    n & l = L \\
    N & \text{otherwise}
  \end{cases}
\end{align*}

$K \equiv M^{-1}$ encodes the prior spatial covariance in synapse strengths: a synapse is defined by its pre-synaptic neuron (at location $\mathbf{x}_a$) and its post-synaptic location (at location $\mathbf{x}_i$) and has correlation with other synapses located by the tuple $(j,b)$; $\mathbb{E}[W_{i,a} W_{j,b}] = K_{i,a,j,b}$. We recover the standard case of $L_2$ regularisation when $K = M = I$, which simply says that our prior is that all synapses are uncorrelated. We can achieve spatial interactions within layers, but not across, by choosing a more general factorisation $K^{(l)}_{i,a,j,b} = K^{(l),\text{post}}_{i,j} K^{(l),\text{pre}}_{a,b}$.
Similarly, though we could choose a more general term for the firing rate regularisation, we will stick with the one that allows only within-layer interactions:
\begin{align*}
  \mathcal{P}_{\text{rates}} = 
    \frac12 \sum_{l=1}^{L} \sum_{\mu=1}^P \sum_{i,j=1}^N \phi^{(l)\mu}_{i} G^{(l)}_{i,j} \phi^{(l)\mu}_{j}
\end{align*}

\subsection{Partition Function}

We then turn our focus to the partition function of \eqref{eq:theta-distr}:
\begin{align*}
  \mathcal Z &=
    \int \mathcal{D} \Theta \mathcal{D} \xi \mathcal{D} \hat \xi
    \exp \Bigg \{
      -i \hat \xi^{(l)\mu}_{i} \left[
        \xi^{(l)\mu}_{i} - \fix{\frac{1}{\sigma^{(l-1)}}} W^{(l-1)}_{i,a} \varphi^{(l-1)\mu}_{a} 
      \right]
      - \frac{\beta N}{2 P n} ( f^{\mu}_{\alpha} - y^{\mu}_{\alpha})^2
      - \mathcal{P}_{\text{weight}}(\Theta)
      - \mathcal{P}_{\text{rates}}(\Phi)
    \Bigg \} 
\end{align*}
where for brevity we have introduced
\begin{align*}
  \xi^{(l)\mu}_{i} = \begin{cases}
    h^{(l)\mu}_{i} & l \in 1, \dots, L \\
    f^{\mu}_i   & l = L+1
  \end{cases}
&&
  \varphi^{(l)\mu}_{i} = \begin{cases}
    x^{\mu}_i         & l =0 \\
    \phi(h^{(l)\mu}_{i}) & l \in 1, \dots, L
  \end{cases}
&& 
  \sigma^{(l)} = \begin{cases}
    \sqrt{m} & l = 0 \\
    N        & l = L \\
    \sqrt{N} & \text{otherwise}
  \end{cases}
\end{align*}
\fixnote{The constraint carried $1/\sigma^{(l)}$ next to $W^{(l-1)}$, but $\sigma$ is indexed by the weight layer (the cases above give $\sigma^{(0)} = \sqrt m$ for $W^{(0)}$ and $\sigma^{(L)} = N$ for the readout), so as written the first layer would be normalised by $\sqrt N$ and the readout by $\sqrt N$. The weight integral below already uses the matching pair $W^{(l)}/\sigma^{(l)}$, so nothing downstream changes; the constraint now reads $W^{(l-1)}/\sigma^{(l-1)}$.}

We begin with the integral over $\Theta$, in effect averaging over the prior \eqref{eq:prior-weights}.
\begin{align*}
  & \int \mathcal D \Theta \exp \left(
    -\frac12 W^{(l)}_{i,a} M^{(l)}_{i,a,j,b} W^{(l)}_{j,b}
    + W^{(l)}_{i,a} \left[
      \frac{i}{\sigma^{(l)}} \hat \xi^{(l+1)\mu}_{i} \varphi^{(l)\mu}_{a}
    \right]
  \right) 
  \propto 
  \exp \left(
    -\typo{\frac{1}{2(\sigma^{(l)})^2}} \hat \xi^{(l+1)\mu}_{i} \varphi^{(l)\mu}_{a} K^{(l)}_{i,a,j,b} \varphi^{(l)\nu}_{b} \hat \xi^{(l+1)\nu}_{j}
  \right)
\end{align*}

We introduce the kernels
\begin{align}
  \Phi^{(0)\mu,\nu}_{i,j}
    = \frac1m \sum_{a,b=1}^m x^{\mu}_a K^{(0)}_{i,a,j,b} x^{\nu}_b
  && \text{and} &&
  \Phi^{(l)\mu,\nu}_{i,j} =
    \frac1N \sum_{a,b=1}^N \phi^{(l)\mu}_{a} K^{(l)}_{i,a,j,b} \phi^{(l)\nu}_{b},
    \quad \forall \ 1 \leq l \leq L
  \label{eq:kernel-def}
\end{align}

which leaves us with the factor
\begin{align*}
  \exp \left( 
    -\frac12 \sum_{l=1}^{L}
      \hat h^{(l)\mu}_{i} \Phi^{(l-1)\mu,\nu}_{i,j} \hat h^{(l)\nu}_{j}
    - \frac{1}{2N} 
      \hat f^{\mu}_{\alpha} \Phi^{(L)\mu,\nu}_{\alpha,\beta} \hat f^{\nu}_{\beta}
  \right)
\end{align*}
Note that, unlike the $\Phi^{(l)}$ kernels we are familiar with, \eqref{eq:kernel-def} introduces two additional degrees of freedom above the usual $\mu,\nu$; these are the $i,j$ indices, which take on values in $1, \dots, \text{Post}(l)$. Let's take a second to examine how this changes the semantics of the $\Phi^{(l)}$ kernels. Suppose we choose a factorisation of $K$ as discussed above, $K_{i,a,j,b} = K^{\text{post}}_{i,j} K^{\text{pre}}_{a,b}$, and additionally, choose it to be diagonalisable in the basis $V$,  $K^{\text{pre}} = V \Lambda^{\text{pre}} V^{*}$. If we now project our firing rates into this basis, $\tilde \phi^{\mu} = V^* \phi^{\mu}$, we have
\begin{align}
  \Phi^{\mu,\nu}_{i,j} = 
  K^{\text{post}}_{i,j} 
  \underbrace{
    \frac1N \sum_q \lambda^{\text{pre}}_q \fix{\tilde \phi^{\mu *}_{q}} \tilde \phi^{\nu}_q
  }_{
    C^{\mu,\nu}
  }
  = 
  K^{\text{post}}_{i,j} C^{\mu,\nu}
  \label{eq:phi-diagonalised}
\end{align}
\fixnote{With $K^{\text{pre}} = V \Lambda^{\text{pre}} V^{*}$ and real rates, $\phi^{\mu\top} K^{\text{pre}} \phi^{\nu} = \sum_q \lambda^{\text{pre}}_q \tilde\phi^{\mu *}_q \tilde\phi^{\nu}_q$; the conjugate was missing, which matters for a complex eigenbasis such as the Fourier basis of Section~\ref{sec:rg} (without it $C$ is not Hermitian). The equation is illustrative and is not used downstream; the conjugate has been added.}
As we will see later, $\Phi^{(l)}$ shows up as the input to layer $l+1$; \eqref{eq:phi-diagonalised} is telling us that the prior we chose for the post-synaptic weights \typo{constrains} how information is sent into that layer \textit{over space} (i.e. carrying the $(i,j)$ index), while the prior we chose for the pre-synaptic weights constrains how information is sent into that layer \textit{over data} (i.e. carrying the $(\mu,\nu)$ index).

To enforce these constraints, we \typo{multiply} $\mathcal Z$ by unity:
\begin{align*}
  1 
  = \int \mathcal D \Phi \int \mathcal D \hat \Phi \exp \left(
    -i \frac{N}{2} \hat \Phi^{(l)\mu,\nu}_{i,j} \left[
      \Phi^{(l)\mu,\nu}_{i,j} - \frac1N \sum_{a,b=1}^N \phi^{(l)\mu}_{a} K^{(l)}_{i,a,j,b} \phi^{(l)\nu}_{b}
    \right]
  \right) 
\end{align*}

Returning to the path integral, we now integrate out $\hat f$
\begin{align*}
  \int \mathcal{D} \hat f \exp \left(
    - \frac{1}{2N} \hat f^{\mu}_{\alpha} \Phi^{(L)\mu,\nu}_{\alpha,\beta} \hat f^{\nu}_{\beta}
    -i \hat f^{\mu}_{\alpha} f^{\mu}_{\alpha} 
  \right)
  \propto 
  \exp \left(
    -\frac{N}{2} f^{\mu}_{\alpha} \big([\Phi^{(L)}]^{-1}\big)^{\mu,\nu}_{\alpha,\beta} f^{\nu}_{\beta}
  \right)
\end{align*}
and subsequently $f$. For this we open up the \typo{quadratic} exponent $||F-Y||^2_2$ and define $\kappa = \frac{\beta}{Pn}$:
\begin{align*}
  & \int \mathcal{D} f \exp \left(
    -\frac{N}{2} f^{\mu}_{\alpha} \left[
      \big([\Phi^{(L)}]^{-1}\big)^{\mu,\nu}_{\alpha,\beta} 
      + \kappa \delta^{\mu,\nu} \delta_{\alpha,\beta}
    \right]
    f^{\nu}_{\beta}
    + N \kappa f^{\mu}_{\alpha} y^{\mu}_{\alpha}
    - \frac{N \kappa}{2} y^{\mu}_{\alpha} (\delta^{\mu,\nu} \delta_{\alpha,\beta}) y^{\nu}_{\beta}
  \right) \\
  & \propto
  \exp \left(
    \frac{N \kappa^2}{2} y^{\mu}_{\alpha} \left[
      (
        [\Phi^{(L)}]^{-1} 
        + \kappa I
      )^{-1}
      - \frac{1}{\kappa} I
    \right]^{\mu,\nu}_{\alpha,\beta} 
    y^{\nu}_{\beta}
  \right)
  =
  \exp \left(
    -\frac{N}{2} y^{\mu}_{\alpha} \left[
      (\Phi^{(L)}
      + \kappa^{-1} I)^{-1}
    \right]^{\mu,\nu}_{\alpha,\beta} 
    y^{\nu}_{\beta}
  \right)
\end{align*}
where the last step follows from the identity $c^2 (A^{-1} + cI)^{-1} - c I = -(A + c^{-1} I)^{-1}$, which is easily proven with algebraic temerity.

The remaining partition function is then
\begin{align*}
  \mathcal Z &=
    \int \mathcal{D} \Phi \mathcal{D} \hat \Phi
    \exp \Bigg \{
     -N \left[
      \frac12 y^{\mu}_{\alpha} \left[ (\Phi^{(L)} + \kappa^{-1} I )^{-1} \right]^{\mu,\nu}_{\alpha,\beta} y^{\nu}_{\beta}
      + \frac{i}{2} \hat \Phi^{(l)\mu,\nu}_{i,j} \Phi^{(l)\mu,\nu}_{i,j}
     \right]
    \Bigg \} \\
    & \qquad \qquad
    \int \mathcal{D} h \mathcal{D} \hat h
    \exp \Bigg \{
      - i \hat h^{(l)\mu}_{i} h^{(l)\mu}_{i} 
      - \frac12 \hat h^{(l)\mu}_{i} \Phi^{(l-1)\mu,\nu}_{i,j} \hat h^{(l)\nu}_{j}
      \fix{{} + \frac12 \phi^{(l)\mu}_{i} \left( i \hat \Phi^{(l)\mu,\nu}_{\alpha,\beta} K^{(l)}_{\alpha,i,\beta,j} - G^{(l)}_{i,j} \delta^{\mu,\nu}  \right) \phi^{(l)\nu}_{j}}
    \Bigg \}
\end{align*}
\fixnote{The rate term had the wrong sign on $i\hat\Phi K$: inserting unity contributes $+i\frac N2 \hat\Phi \cdot \frac1N \phi K \phi$ and the rate penalty $-\frac12 \phi G \phi$, so the exponent is $+\frac12 \phi(i\hat\Phi K - G)\phi$. Kept, the old sign would give $\Phi = -\frac1N\langle\phi K \phi\rangle$ at the saddle; \eqref{eq:single-site} below already carries the correct sign, so nothing downstream changes.}

\paragraph{Input Noise}

Until now I have suppressed the $\eta$ that were introduced \typo{in} \eqref{eq:ff-network}. However, we can here take an expectation over them to get
\begin{align}
  \int \mathcal D \eta \exp \left(-\frac{1}{2\sigma^2} \eta^{\mu}_i [\delta^{\mu,\nu} \delta_{i,j}] \eta^{\nu}_{\typo{j}} + i \eta^{\mu}_i \hat h^{(1)\mu}_i \right)
  \propto
  \exp \left(-\frac{\sigma^2}{2} [\hat h^{(1)\mu}_i]^2 \right)
  \label{eq:single-site-reg}
\end{align}
This term will regularise the input covariance to the first hidden layer. In general, I will have my cake and eat it by generally considering the limit $\sigma \rightarrow 0^+$ and everything is nicely invertible, and so suppress it in the notation for the single-site below.

\subsection{Single-Site}

We can now introduce the independent (conditioned on the kernels $\Phi$ and $\hat \Phi$) sites for each layer. Note that although these single-sites should in general carry the layer-specific $K$ and $G$ that is germane, I will drop this index:
\begin{align}
  \mathcal W[\Phi, \hat \Phi] &= 
  \int \mathcal{D} h \mathcal{D} \hat h
    \exp \Bigg \{
      - i \hat h^{\mu}_{i} h^{\mu}_i 
      - \frac12 \hat h^{\mu}_i \Phi^{\mu,\nu}_{i,j} \hat h^{\nu}_j
      + \frac12 \phi^{\mu}_i \left( i \hat \Phi^{\mu,\nu}_{\alpha,\beta} K_{\alpha,i,\beta,j} - G_{i,j} \delta^{\mu,\nu}  \right) \phi^{\nu}_j
    \Bigg \} \\
  &= 
  \int \mathcal{D} h \exp \Bigg \{
    - \frac12 h^{\mu}_i (\Phi^{-1})^{\mu,\nu}_{i,j} h^{\nu}_j 
    + \frac12 \phi^{\mu}_i \left( i \hat \Phi^{\mu,\nu}_{\alpha,\beta} K_{\alpha,i,\beta,j} - G_{i,j} \delta^{\mu,\nu}  \right) \phi^{\nu}_j
  \Bigg \} \label{eq:single-site}
\end{align}
where in going from the first to the second line we have completed the gaussian integral over $\hat h$ (which is convergent in the first layer with the inclusion of \eqref{eq:single-site-reg}). Our goal will be to find the free energy density of these single-site distributions
\begin{align*}
  \mathcal W[\Phi, \hat \Phi] \asymp e^{-N \mathcal F[\Phi, \hat \Phi]}
\end{align*}
such that we can write
\begin{align*}
  \mathcal Z &=
    \int \mathcal{D} \Phi \mathcal{D} \hat \Phi
    \exp \Bigg \{
      -N \left[
      \frac12 y^{\mu}_{\alpha} \left[ (\Phi^{(L)} + \kappa^{-1} I )^{-1} \right]^{\mu,\nu}_{\alpha,\beta} y^{\nu}_{\beta}
      + \frac{i}{2} \hat \Phi^{(l)\mu,\nu}_{i,j} \Phi^{(l)\mu,\nu}_{i,j}
      + \sum_{l=1}^L \mathcal F[\Phi^{(l-1)}, \hat \Phi^{(l)}]
     \right]
    \Bigg \}
\end{align*}
and approximate the full partition function by a saddle-point approximation.
In general, we can construct the free energy density as follows:
\begin{align*}
  \mathcal W[\Phi, \hat \Phi]
  &= 
  \major{\int \mathcal{D} h}
  \prod_{i=1}^N
  \underbrace{
    \major{\exp} \Big \{
      - \sum_{j=1}^N \frac12 h^{\mu}_i (\Phi^{-1})^{\mu,\nu}_{i,j} h^{\nu}_j 
      \fix{{} + \frac12 \phi^{\mu}_i \left( i \hat \Phi^{\mu,\nu}_{\alpha,\beta} K_{\alpha,i,\beta,j} - G_{i,j} \delta^{\mu,\nu}  \right) \phi^{\nu}_j}
    \Big \}
  }_{\major{w_i[h]}}
\end{align*}
\fixnote{Same sign slip on the $i \hat \Phi K$ term as in the remaining partition function above; corrected to match \eqref{eq:single-site}.}
\major{The integrand factorises over space, but the integral does not: each $w_i[h]$ depends on the preactivations $h_j$ of every site $j$ coupled to $i$ through $\Phi^{-1}$ or $K$, so $\mathcal W$ is not a product of single-neuron integrals and $\mathcal F$ is not an average of single-neuron free energies.}
\begin{align*}
  -\mathcal F[\Phi, \hat \Phi] 
  =
  \frac1N \log \mathcal W[\Phi, \hat \Phi] 
  =
  \major{\frac1N \log \int \mathcal D h \prod_{i=1}^N w_i[h]}
\end{align*}
\major{A genuine average appears only after a change of basis that decouples the sites, e.g.\ over Fourier modes for a translation-invariant linear layer (Section~\ref{sec:rg}).}\majornote{\detokenize{setup_site_factorisation.tex}}
This is for now a conceptual \typo{aid}, rather than a genuine result: $\major{w_i}$ still includes an extensive sum in its exponent, and so we have no reason to believe \major{$\mathcal F$} is actually $O(1)$. More work will be required to find a convergent density.

\subsection{Saddle-Point}

As $N \rightarrow \infty$, we can approximate $\mathcal Z$ by the saddle-point method. Consequently, we find the stationary points with respect to the kernels and their conjugates. \major{At the saddle point $\hat \Phi$ is purely imaginary (the readout condition below fixes $i \hat \Phi^{(L)}$ to be real), so we first rotate the contour, $\hat \Phi \rightarrow -i \hat \Phi$: the term $\frac{i}{2} \hat \Phi \Phi$ in $\mathcal Z$ becomes $\frac12 \hat \Phi \Phi$ and $i \hat \Phi K$ in $\mathcal W$ becomes $\hat \Phi K$, so that $\mathcal S$ and $\mathcal W$ are real. Every $\hat \Phi$ from here on, and in Section~\ref{sec:rg}, is the rotated one.}\majornote{\detokenize{setup_conjugate_rotation.tex}} Defining the action
\begin{align*}
  \mathcal Z &= \int \mathcal D \Phi \mathcal D \hat \Phi \exp \Big( - N \mathcal S[\Phi, \hat \Phi] \Big) \\
  \mathcal S &= 
    \frac12 y^{\mu}_{\alpha} 
      \left. (\Phi^{(L)} + \kappa^{-1} I )^{-1} \right|^{\mu,\nu}_{\alpha,\beta} y
    ^{\nu}_{\beta}
    + \frac{1}{2} \hat \Phi^{(l)\mu,\nu}_{i,j} \Phi^{(l)\mu,\nu}_{i,j}
    \major{{} - \frac1N \sum_{l=1}^L \log \mathcal W[\Phi^{(l-1)}, \hat \Phi^{(l)}]}
\end{align*}
we obtain the self-consistency relation
\begin{align*}
  \frac{\partial \mathcal S}{\partial \hat \Phi^{(l)}} = 0
  \implies \Phi^{(l)\mu,\nu}_{i,j}
    = \major{\frac1N \sum_{a,b=1}^N K^{(l)}_{i,a,j,b} \left \langle \phi^{\mu}_a \phi^{\nu}_b \right \rangle^{(l)}}
\end{align*}
where \typo{we} have defined $\major{\langle \cdot \rangle^{(l)}}$ as the expectation over the \typo{neuron} response distribution \eqref{eq:single-site} for layer $l$ (i.e. with input drive $\Phi^{(l-1)}$ and tilting $\hat \Phi^{(l)}$), \major{taken jointly over all neurons of the layer, which are coupled}.\majornote{\detokenize{setup_site_factorisation.tex}} For the tilting we obtain for $l \in 1, \dots, L-1$.
\begin{align}
  \frac{\partial \mathcal S}{\partial \Phi^{(l)}} = 0
  \implies \hat \Phi^{(l)\mu,\nu}_{i,j} 
  = \frac1N \sum_{k=1}^N \langle \hat h^{\mu}_i \hat h^{\nu}_j \rangle_{k}^{(l+1)}
  = \frac1N \sum_{k=1}^N \left \langle
    \left. \frac{\partial h}{\partial \eta} \right|^{\mu}_i 
    \left. \frac{\partial h}{\partial \eta} \right|^{\nu}_j 
  \right \rangle_{k}^{(l+1)}
  \label{eq:saddle-phi-hat}
\end{align}
where $\eta$ is the input current to neuron $h$. Alternatively, using the form \eqref{eq:single-site} with $\hat h$ integrated out, we obtain
\begin{align}
  \hat \Phi^{(l)} = [\Phi^{(l)}]^{-1} \left( \frac1N \sum_{k=1}^N \langle h \otimes h \rangle_{k}^{(l+1)} \right) [\Phi^{(l)}]^{-1}
  \label{eq:backward-pass}
\end{align}
For $l=L$, $\Phi^{(L)}$ does not appear in any single site, only the tilting from the readout, giving us
\begin{align*}
  \hat \Phi^{(L)} = 
    (\Phi^{(L)} + \kappa^{-1} I)^{-1}
    \mathbb Y
    (\Phi^{(L)} + \kappa^{-1} I)^{-1}
\end{align*}
where we have defined $\mathbb Y^{\mu,\nu}_{\alpha,\beta} = y^{\mu}_{\alpha} y^{\nu}_{\beta}$.

\subsection{Predictive Distribution}

The last element of the standard mean-field treatment is to find the distribution over test outputs $f^{\star}_{\alpha}$ corresponding to test inputs $\typo{x^{\star}_a}$. This has the standard form of a Gaussian process regression:
\begin{align*}
  f^{\star}_{\alpha} | \Phi^{(L)} \sim
    \mathcal N (\mu^{\star}_{\alpha}, \frac1N \Sigma^{\star}_{\alpha,\beta})
\end{align*}
with moments
\begin{align*}
  \mu^{\star}_{\alpha} & = 
    \Phi^{(L)\star,\mu}_{\alpha,\fix{\alpha'}}
    \left. \fix{(\Phi^{(L)} + \kappa^{-1} I)^{-1}} \right|^{\mu,\nu}_{\fix{\alpha',\beta'}}
    y^{\nu}_{\fix{\beta'}}
  \\ 
  \text{ and }
  \Sigma^{\star}_{\alpha,\beta} & =
    \Phi^{(L)\star,\star}_{\alpha,\beta} 
    - \Phi^{(L)\star,\mu}_{\alpha,\fix{\alpha'}}
    \left. \fix{(\Phi^{(L)} + \kappa^{-1} I)^{-1}} \right|^{\mu,\nu}_{\fix{\alpha',\beta'}}
    \Phi^{(L)\nu,\star}_{\fix{\beta'},\beta}
\end{align*}
\fixnote{The regularised kernel must be inverted, as in any GP posterior: with prior $f \sim \mathcal N(0, \Phi^{(L)}/N)$ and likelihood variance $1/(N\kappa)$ the mean is $\Phi^{(L)\star\mu}[(\Phi^{(L)}/N + I/(N\kappa))^{-1}]^{\mu\nu} y^\nu/N$, which is the expression above. The summed indices live in output space, so they are now output indices $\alpha',\beta'$ (they were the input labels $a,b$, and the last factor carried $a$ where the contraction needs the second index). Nothing in the saddle-point equations depends on this side result.}
as $N \rightarrow \infty$, this concentrates at $\mathbf{\bar f}^{\star}$

\end{document}
